\section{Discussion and Conclusion}
\label{sec:discussion}

\paragraph{Practical implications.}
Our findings provide clear guidance for practitioners: when LP performance on sparse or tail regions matters, standard graph SSL (global CL or even simple augmentation) is a robust default. There is no need to design degree-aware or selective CL methods---the benefit is inherently global and disproportionately helps tail nodes. On large, dense graphs, augmentation alone suffices, avoiding the computational overhead of contrastive objectives.

\paragraph{Limitations.}
Our study has several important limitations. First, we cover three message-passing backbones (GCN, GraphSAGE, GAT); the phenomenon may differ for graph transformers, subgraph-based LP methods, or other encoder families. Second, while our partial correlation analysis (Table~\ref{tab:partial}) shows that degree independently predicts SSL gain after controlling for clustering and neighborhood diversity, degree remains entangled with other unmeasured structural properties---fully disentangling these factors requires further investigation. Third, the controlled sparsity experiment uses both uniform and degree-preserving edge dropping, which produce consistent results but remain artificial interventions that may not fully capture naturally sparse graph structures. Fourth, our inductive evaluation (Table~\ref{tab:inductive}) confirms CL benefits generalize to unseen nodes, though the setup uses feature-based generalization and does not test distribution-shifted or temporal settings. Fifth, our benchmarks cover citation, coauthor, co-purchase, and collaboration graphs; extension to heterogeneous and dynamic graphs remains open.

\paragraph{Conclusion.}
We presented a systematic degree-stratified analysis of graph self-supervision for link prediction. Our central finding---that SSL gains are monotonically concentrated on tail nodes---is robust across 6 datasets, 2 backbones, and multiple training methods. The benefit decomposes into augmentation regularization and the contrastive objective, with their relative importance depending on graph sparsity. Selective CL targeting does not improve upon global CL, and non-CL alternatives provide substantially smaller gains. These results challenge the assumption of uniform SSL benefits and highlight the value of degree-stratified evaluation as a standard practice in graph LP research.
